\chapter{La particule comme section matérielle persistante : histoire, représentation, ruban, composition et critère d’identité}
\label{chap:particula-seccion-material-persistente}

\section{Définition de la particule comme classe compatible de sections matérielles persistantes sous raffinement}
Une particule est une classe compatible de sections matérielles persistantes de l’extension généalogique, dotée d’une histoire observable, d’une représentation, d’une orientation et d’un critère d’identité sous raffinement. Le parcours est son histoire, non la particule elle-même ; la cellule est une coordonnée de l’atlas, non un corps qui se déplace ; la masse est une observable spectrale, non sa définition ontologique.

\section{Classificateur des espèces matérielles par fibre, représentation, charge, spin, saveur et génération}
L’espèce est typée par la charge, le spin, la couleur, la saveur, la génération, la représentation et l’orientation. La charge est un caractère orienté ; le spin est une donnée de représentation ; la couleur est une fibre ternaire résiduelle ; saveur et génération relèvent de la multisection et de sa profondeur. Aucune de ces données ne se réduit à la valeur de masse.

\section{Antiparticule, ruban et virtualité}
L’involution \(C_M=-\operatorname{rev}\) induit l’antisection et la conjugaison des feuillets. Particule et antiparticule peuvent être deux sections d’un même revêtement de Möbius lorsque le lacet, l’holonomie et la représentation correspondants sont fixés. Le retour étendu de 216 pas ne se décompose pas en deux trajectoires indépendantes de 108. Une particule virtuelle est une section intermédiaire non asymptotique à horizon fini de prolongement.

\section{Persistance de trajectoire, multisections et conjugaison de Möbius}

La persistance se formule comme section d’un système inverse ; l’espèce interne comme quotient orbital d’une multisection ; et la conjugaison comme antisection à holonomie déclarée. Le développement complet conserve la stratification extension--parcours--registre--signature, le lecteur temporel et la distinction entre les deux occurrences de \(2\pi/4\pi\).

\begingroup
\let\chapter\section
\let\section\subsection
\let\subsection\subsubsection
\let\subsubsection\paragraph
\let\clearpage\relax
\let\cleardoublepage\relax
\section{Persistance de voie, multisections d'espèce et conjugaison de Möbius}
\label{sec:persistencia-ruta-mobius-familias}

Cette section précise une distinction qui, formulée de manière trop
abrégée, inverse la logique de la construction. La voie n'est pas une donnée qu'une
particule choisit une fois sa masse connue. Ce n'est pas non plus un dessin auxiliaire
sans contenu causal. C'est l'histoire observable et orientée d'une extension
survivante. L'extension conserve, en outre, la profondeur compatible, le
registre, les retenues et la mémoire qui ne tiennent pas dans une coupe finie. En ce
sens précis,
\[
 \boxed{\text{une particule HMT est la persistance compatible d'une route
 enrichie}.}
\]
La phrase et la définition par extensions expriment le même objet à deux
niveaux : la voie est son histoire visible ; l'extension est la tour complète de
cette histoire avec tous ses prolongements compatibles.

\subsection{Persistance comme section d'un système inverse}

Soit \((\mathcal S_m,\rho_{m+1,m})_{m\geq0}\) le système d'états qui
survivent jusqu'à la profondeur \(m\), et soit
\[
 x=(x_m)_{m\geq0}\in\Omega_{\rm HMT}^{\rm surv}
 =\varprojlim_m\mathcal S_m.
\]
Les ombres de voie satisfont le carré de compatibilité
\[
 \operatorname{tr}_{m+1,m}\!\left(\operatorname{sh}_{m+1}(x_{m+1})\right)
 =\operatorname{sh}_{m}\!\left(\rho_{m+1,m}(x_{m+1})\right)
 =\operatorname{sh}_{m}(x_m).
\]
Par conséquent
\[
 \gamma_x=\bigl(\gamma_x^{[m]}\bigr)_{m\geq0},
 \qquad \gamma_x^{[m]}=\operatorname{sh}_m(x_m),
\]
est une histoire compatible, non une séquence ajustée rétrospectivement.

Sur chaque coupe \(\gamma_x^{[m]}\), on dispose du fibré discret
\(\mathcal B_{\gamma_x^{[m]}}\). Une section persistante est une famille
\[
 s=(s_m)_{m\geq0},
 \qquad
 s_m\in\Gamma\!\left(\mathcal B_{\gamma_x^{[m]}}\right),
 \qquad
 \operatorname{res}_{m+1,m}(s_{m+1})=s_m.
\]
Cette équation est le contenu mathématique de la persistance : chaque lecture
finie se prolonge sans contredire les lectures précédentes. Le tuple de
particule du \cref{chap:particula} ajoute l'occupation,
la représentation et la propagation, mais n'altère pas cette provenance.

\subsection{Espèce interne comme quotient orbital de la multisection}

Dans l'atlas nonadique, chaque cellule orientée \(c=(r,c')\in\mathcal C_{81}\)
porte une famille de voie
\[
 \mathcal R(c)=
 (n_A,s_C,k_{\Delta_4},\nu_{120},\nu_{270},\tau_{12}).
\]
L'image de \(\mathcal R\) contient cinquante-six valeurs. Si l'on fixe la
carte et son ordre, le rang \(\kappa\in\{0,1,2,3,4\}\) au sein de chaque fibre
produit le relèvement injectif
\[
 c\longmapsto \bigl(\mathcal R(c),\kappa(c)\bigr).
\]
L'alphabet à cinq états est minimal, car il existe deux fibres de
cardinal cinq. Cette sélection est interne, prospective et aveugle aux masses et aux
noms de particules.

Le recensement complet qui soutient cette affirmation est
\[
 81=45_{\rm leptones}+36_{\rm quarks}
   =25_{\ell^-}+20_{\nu}+20_{d\text{-type}}+16_{u\text{-type}}.
\]
Les cinquante-six fibres de voie possèdent l'histogramme
\[
 41\times1+10\times2+2\times3+1\times4+2\times5=81.
\]
Par conséquent, \(56\) compte les signatures de voie, tandis que \(81\) compte les cellules
orientées ; la mémoire \(\kappa\) distingue les cellules qu'une même signature ne
sépare pas.

Pour classer les espèces internes, on n'inverse pas cette application. On prend le
quotient
\[
 q:\mathcal C_{81}\longrightarrow\Sigma_{13},
 \qquad
 q(c)=\bigl(\operatorname{familia}(c),G(c)\bigr),
\]
dont l'image comprend treize types. La sortie associée à un type \(y\) est la
multisection finie de l'atlas
\[
 \boxed{
 \mathfrak S_y^{\rm atlas}=
 \{(\mathcal R(c),\kappa(c)):q(c)=y\}.}
\]
Ses cardinaux sont
\[
 1,8,16;\qquad 2,4,6,8;\qquad 2,4,6,8;\qquad4,12,
\]
pour les leptons chargés, les neutrinos et les deux types de quark,
respectivement. L'espèce est le type commun de cette fibre ; une particule
concrète est une section persistante réalisée en son sein. Exiger que le
nom d'une espèce choisisse à lui seul une cellule détruit précisément les
degrés d'orientation et de conjugaison que la fibre conserve.

Le relèvement aux extensions est un objet typé différent. Si
\(\pi_{81}:\Omega_{\rm HMT}^{\rm surv}\to\mathcal C_{81}\) désigne le lecteur
cellulaire de la voie enrichie déclarée, alors
\[
 \widetilde{\mathfrak S}_y
 =\{x\in\Omega_{\rm HMT}^{\rm surv}:q(\pi_{81}x)=y\}
\]
est la multisection enrichie correspondante. Le certificat fini
démontre le recensement et l'obstruction dans
\(\mathfrak S_y^{\rm atlas}\) ; son relèvement est relatif au lecteur
\(\pi_{81}\), non une identification automatique entre une cellule et une
histoire profonde.

\begin{theorem}[Obstruction au sélecteur ponctuel et exception électronique]
Soit \(\rho_{\rm c}(r,c')=(10-r,10-c')\) l'inversion centrale. Les treize fibres de
\(q\) sont \(\rho_{\rm c}\)-stables. La fibre de lepton chargé de niveau zéro possède
l'unique point fixe, \((5,5)\). Aucune des douze autres fibres n'admet
un sélecteur ponctuel \(\rho_{\rm c}\)-équivariant si \(\rho_{\rm c}\) agit trivialement
sur le type grossier.
\end{theorem}
\begin{proof}
L'équation \(\rho_{\rm c}(r,c')=(r,c')\) équivaut à \(r=c'=5\). Le recensement de l'atlas
situe cette cellule dans la fibre \((\ell^-,G0)\), qui a pour cardinal un. Soit
maintenant \(y\neq(\ell^-,G0)\) et supposons qu'il existe une section ponctuelle
équivariante \(a:\Sigma_{13}\to\mathcal C_{81}\). Comme \(\rho_{\rm c} y=y\),
l'équivariance exigerait
\[
 a(y)=a(\rho_{\rm c} y)=\rho_{\rm c} a(y).
\]
Alors \(a(y)\) devrait être un point fixe de \(\rho_{\rm c}\), mais l'unique
point fixe appartient à la fibre électronique. C'est impossible. Par conséquent,
dans les douze autres fibres, la sortie naturelle est une \(\rho_{\rm c}\)-orbite ou une
multisection ; choisir un représentant requiert une orientation ou une
représentation supplémentaire.
\end{proof}

L'électron est exceptionnel en trois sens compatibles : \((5,5)\) est le
centre géométrique unique, l'unique cellule de niveau zéro et le centre par rapport
auquel se forme le réseau brut de différences
\[
 \Lambda_{\rm raw}^{\Delta}
 =\bigl\langle A_{\rm raw}(c)-A_{\rm raw}(5,5):
 c\in\mathcal C_{81}\bigr\rangle_{\mathbb Z}.
\]
La signature brute de cette cellule n'est pas nulle :
\[
 A_{\rm raw}(5,5)=(24,0,12,0,-12),
 \qquad \tau(5,5)=\texttt{+++---+++---}.
\]
L'origine \(v_e=(0,0,0,0,0)\) du réseau d'action est une convention
relative à l'électron et non la signature brute de la cellule centrale.

L'invariant algébrique électronique est évalué au moyen de
\[
 \mathfrak e_0=
 \frac{\sqrt3}{4}
 \left(e^{\varphi/\pi^2}-22\alpha_{\rm HMT}^{3}\right)
 (1+15\Delta_4),
 \qquad
 \Delta_4=
 \frac{\pi-e\log\pi}{270}\left(1+\frac{\pi}{729}\right),
\]
\[
 \mathfrak e_0
 =0.510998922488066072509870877191349\ldots.
\]
L'expression est adimensionnelle ; la carte
\(\varsigma_{u_E}:y\mapsto y\,u_E\) introduit ensuite une unité d'énergie.
Le centre \((5,5)\) et son unicité sont des résultats finis exacts. La sélection
des coefficients \(22\) et \(15=3\cdot5\) se déduit, avec la structure de
départ explicite, au moyen du complément
des cinq régions de \(\pi\) dans \(K_e=\mathbb Z/9\mathbb Z\times\mathbb F_3\) :
\(-22=-|K_e\setminus\iota(R_\pi)|\) et
\(+15=|R_\pi\times\mathbb F_3|\). La signature brute centrale conserve un type
distinct et agit comme contrôle indépendant de la sélection basale.

La cellule \((5,5)\) ne doit pas être confondue avec la signature énergétique
\(555555\) : cette dernière appartient au représentant transversal d'une des
cinq régions de \(\pi\). Ce sont des objets distincts qui partagent le chiffre
central.

Dans le secteur des quarks, le caractère résiduel
\[
 \operatorname{col}(r,c')=r-c'\pmod3
\]
raffine la multisection en trois classes de douze cellules et satisfait
\(\operatorname{col}(\rho_{\rm c} c)=-\operatorname{col}(c)\). Cela construit le
tritique interne de couleur. La représentation physique de \(SU(3)\) est un
lecteur ultérieur ; elle n'est pas nécessaire pour démontrer la répartition finie.

\subsection{Stratification d'extension, de voie, de registre et de signature}

Les applications \(q\), \(\mathcal R\), la projection cellulaire et le lecteur de
réalisation distinguent quatre classifications liées et typées :
\begin{enumerate}[label=(\roman*)]
 \item les treize fibres internes famille--niveau de \(q\) ;
 \item les cinquante-six familles de voie de \(\mathcal R\) ;
 \item les quatre-vingt-une positions et leurs trois cent vingt-quatre
       orientations ;
 \item les réalisations physiques individualisées par représentation et codomaine.
\end{enumerate}
Les quatre couches sont recensées avant la confrontation métrologique. Pour chaque classe,
la comparaison requiert le tuple
\[
(\text{position},\text{fibre},\text{routes},\text{lecteurs},
\text{registre},\text{signature},\text{caractère},\text{tours},
\text{carte dimensionnelle}).
\]
Une reconstruction sur un sous-ensemble démontre seulement ce sous-cas et ne
remplace aucune des composantes du tuple.
Étant donnés \(A\), \(C^*\) et les canaux \(q_\pm\), la tour globale appartient au
semi-groupe infini
\[
 s_n=q_-^n-q_+^n,
 \qquad
 n\in\langle90,120\rangle
 =\{90,120\}\cup\{30r:r\geq6\}.
\]
Les hauteurs \(90,120,180,210,240,270,360,420\) sont des représentants finis du
semi-groupe ; la récurrence précédente prolonge la tour complète.
et le caractère adimensionnel d'une signature est
\[
 \chi_{\rm mass}(v)=\exp\!\left(
 n_AA_{\rm rad}+\frac{s_C}{6}C^*_{\rm rad}
 +k\Delta_4+\nu_{120}s_{120}+\nu_{270}(135\Delta_4^2)
 \right).
\]
Les moments sont globaux : les espèces les pondèrent au moyen de leur voie, de leur
registre et de leur opérateur de fibre. Les lecteurs \(K_D\) et \(K_\Omega\)
déterminent le raffinement bidirectionnel avant la confrontation, tandis que la
table intégrale des espèces publie séparément la généalogie interne, la
carte dimensionnelle et la comparaison externe.

\subsection{Conjugaison, anti-section et ruban de Möbius}

Sur un mot dodécaphasique, on définit l'involution exacte
\[
 C_M(\tau)_m=-\tau_{11-m},
 \qquad C_M^2=I.
\]
Les fonctionnelles linéaires à poids symétriques sont impaires sous \(C_M\),
tandis que les corrélations quadratiques compatibles avec le renversement sont
paires. Ainsi, l'orientation et la charge linéaire vivent dans le secteur impair,
tandis que la mémoire quadratique vit dans le secteur pair. Cela ne démontre pas que
le caractère exponentiel complet de masse soit pair : si
\(\chi(\tau)=e^{\ell(\tau)}\) et que \(\ell\) est impair, alors
\(\chi(C_M\tau)=\chi(\tau)^{-1}\).

L'involution précédente agit exactement sur les mots. Pour parler
d'anti-section, on fixe en outre un relèvement \(\widetilde C_M\) au fibré des
sections qui préserve l'admissibilité, le registre et les troncatures :
\[
 \operatorname{res}_{m+1,m}\widetilde C_M
 =\widetilde C_M\operatorname{res}_{m+1,m},
 \qquad
 \operatorname{sh}(\widetilde C_Mx)=C_M\operatorname{sh}(x).
\]
Ce n'est que sous ce relèvement que l'anti-section de \(s\) est
\(\widetilde C_Ms\) sur la voie inversée.

Le mot ``Möbius'' possède ici une formulation relative précise. Soit
\(\ell_{\rm ret}\) une boucle de retour de la voie enrichie qui engendre
\[
 \langle\ell_{\rm ret}\rangle\simeq\pi_1(S^1)\simeq\mathbb Z.
\]
Ce groupe de retour n'est pas la phase finie \(C_9\). Supposons que son
relèvement au
registre possède un caractère de signe
\[
 \varepsilon_x:\langle\ell_{\rm ret}\rangle\longrightarrow\{\pm1\},
 \qquad \varepsilon_x(\ell_{\rm ret})=-1.
\]
Le fibré en intervalles associé est
\[
 \mathcal M_x=
 ([0,1]\times[-1,1])/(0,u)\sim(1,-u).
\]
Sa première classe de Stiefel--Whitney est non triviale ; par conséquent
\(\mathcal M_x\) est un ruban de Möbius. Un tour échange les deux
feuillets locaux et deux tours restaurent l'orientation. Cela n'impose pas que
le registre complet revienne au même état : retour d'orientation et retour
de mémoire demeurent des types différents.

Le théorème est relatif au relèvement, à la boucle de retour et au caractère
\(\varepsilon_x\), tous déclarés. Il ne naît pas de l'observation de deux signes.
La lecture des feuillets requiert, en outre, de distinguer la graduation scindée \(R\)
de la conjugaison \(C\) et d'imposer
\[
 CRC^{-1}=-R,
 \qquad CT(C^*)C^{-1}=T(-C^*).
\]
La lecture particule--antiparticule exige finalement que la représentation
physique inverse les porteurs impairs, préserve les porteurs quadratiques
pairs et entrelace l'évolution et le caractère complet de la manière déclarée
par la représentation. Sous ces prémisses, particule et antiparticule sont
les deux orientations conjuguées d'une même bande de persistance.

\begin{theorem}[Décomposition paire--impaire sur un revêtement double]
\label{pf84:A10-md-mobius}
Soit \(p:\widetilde X\to X\) un revêtement principal double d'involution
\(\iota\). Pour toute \(f\in C^0(\widetilde X,\mathbb R)\), les projecteurs
\[
 f^+=\frac12(f+\iota^*f),\qquad
 f^-=\frac12(f-\iota^*f)
\]
donnent la décomposition \(f=f^++f^-\), avec
\[
 f^+\circ\iota=f^+,\qquad f^-\circ\iota=-f^-.
\]
La partie paire descend à une fonction sur \(X\) ; la partie impaire est une
section du fibré en droites de signe associé. Pour le revêtement d'orientation
du ruban de Möbius, \(w_1\ne0\) fait obstruction à une section globale continue
et non nulle de ce fibré en droites.
\end{theorem}

\begin{proof}
Les identités des deux projecteurs prouvent la somme et les parités. Une
fonction invariante est constante sur les fibres de \(p\), donc
descend ; une fonction anti-invariante satisfait la loi de transition de la
représentation de signe. Dans le ruban de Möbius, le fibré en droites associé est non
trivial et sa première classe de Stiefel--Whitney est non nulle, de sorte que toute
section continue globale possède un zéro. La conclusion exclut une
section globale continue \emph{sans zéros} ; elle ne nie pas l'existence de
sections avec des zéros.
\end{proof}

Ce lemme classique type le secteur pair--impair du revêtement ; il n'identifie pas
à lui seul le renversement de mot, l'holonomie de retour, la
graduation scindée, la conjugaison des feuillets ni la parité d'une signature de
masse. Ces identifications requièrent leurs propres opérateurs d'entrelacement.

\subsection{TRIT algébrique et lecteur temporel}

La réalisation quadratique déclarée associe au signe visible du TRIT un type
algébrique. Pour
\(\bar\tau\in\{+1,0,-1\}\), soit
\[
 \mathcal A_{\bar\tau}=
 \mathbb R[J]/(J^2+\bar\tau).
\]
Alors
\[
 \begin{array}{c|c|c}
 \bar\tau&J^2&\text{régime}\ \\ \hline
 +1&-1&\text{elliptique/complexe}\ \\
 0&0&\text{parabolique/nilpotent}\ \\
 -1&+1&\text{hyperbolique/scindé}.
 \end{array}
\]
Les exponentielles correspondantes sont
\[
 e^{tJ}=\cos t+J\sin t,
 \qquad e^{tJ}=1+tJ,
 \qquad e^{tJ}=\cosh t+J\sinh t,
\]
et dans le secteur scindé apparaissent les projecteurs
\(P_\pm=(1\pm J)/2\). Ces identités se vérifient exactement dans la
réalisation déclarée. La
courbure ne provient pas du signe isolé : elle apparaît lorsque APP apporte des liens et
des deux-cellules et que l'on fixe une connexion. Dans la polarisation mixte somme--produit
publiée, sur les \(81\) plaquettes, la courbure orientée satisfait
\[
 F(S,S)=F(P,P)=0,
 \qquad F(S,P)=+1,
 \qquad F(P,S)=-1\pmod9.
\]
L'identité de Stokes finie correspondante contrôle que le bord de chaque
rectangle reproduit la somme de ses plaquettes.

En ordonnant les transitions, on obtient le lecteur temporel tritique :
\[
 \begin{aligned}
 +1&\longmapsto\text{retour, mémoire et passé},\\
  0&\longmapsto\text{seuil torsionnel et présent},\\
 -1&\longmapsto\text{ouverture bidirectionnelle et futur}.
 \end{aligned}
\]
Son précurseur exact est la différence entre retour de phase et retour
d'état :
\[
 g(k+9)=g(k),
 \qquad B_{k+9}\neq B_k\quad\text{en général}.
\]
Dans la réalisation géométrique correspondante, on parle de secteur
sphérique--volumétrique, de seuil plan avec registre torsionnel et de secteur
hyperbolique--surfacique. Cette correspondance définit un lecteur de mécanique
dimensionnelle soumis à l'application entre le TRIT et la géométrie ; elle ne redéfinit pas le TRIT
algébrique et n'identifie pas à elle seule
un signe à une courbure de l'espace-temps. L'expression ``cristal
temporel'' nomme cette architecture dans la terminologie HMT ; une réalisation
dynamique de cristal temporel exige en outre un hamiltonien ou un protocole
périodique d'excitation,
une réponse sous-harmonique robuste et une stabilité.

\subsection{Support exact de la lecture aller--retour}

L'échange bidirectionnel ne repose pas seulement sur une image verbale. Les
coordonnées angulaires entières
\[
 W_+=6n_A+s_C,
 \qquad W_-=6n_A-s_C
\]
s'obtiennent au moyen de
\[
 \binom{W_+}{W_-}
 =\begin{pmatrix}6&1\\6&-1\end{pmatrix}
 \binom{n_A}{s_C},
 \qquad
 \operatorname{SNF}\!\begin{pmatrix}6&1\\6&-1\end{pmatrix}
 =\operatorname{diag}(1,12).
\]
Le réseau est donc d'indice douze. L'involution
\(s_C\mapsto-s_C\) échange \(W_+\leftrightarrow W_-\) et, de manière
équivalente, \(q_+\leftrightarrow q_-\). Tel est le substrat algébrique exact
de l'aller et du retour. Son interprétation comme propagation avancée et retardée
requiert l'opérateur d'évolution de la représentation physique.

La projection paire employée dans la lecture biangulaire satisfait, pour toute
fonction pour laquelle les expressions sont définies,
\[
 \frac{f(\phi)+f(-\phi)}2=f_{\rm par}(\phi),
 \qquad
 \frac{(\pi-\phi)^2+(\pi+\phi)^2}{2}=\pi^2+\phi^2.
\]
Ce sont des identités exactes de symétrisation. La lecture en termes d'une
théorie physique de l'absorbeur relève du foncteur de propagation, non de
l'identité algébrique isolée.

\subsection{Périodicités spinorielles non équivalentes de \texorpdfstring{\(2\pi\)}{2 pi} et \texorpdfstring{\(4\pi\)}{4 pi}}

Dans le relèvement spinoriel classique,
\[
 U(\theta)=\exp\!\left(-\frac{i\theta}{2}\sigma_z\right),
 \qquad U(2\pi)=-I,
 \qquad U(4\pi)=I.
\]
L'égalité est exacte dans le revêtement double \(SU(2)\to SO(3)\). Son application
à une section HMT requiert de spécifier la représentation spinorielle du
transport. Une fois fixée, elle explique pourquoi un tour change le signe de la
section et deux tours le restaurent. Si l'on adopte les relations CAR dans
l'algèbre d'occupation, l'exclusion s'en déduit ; ni CAR ni Pauli ne
s'infèrent du signe de \(2\pi\) sans cette structure supplémentaire.

L'interprétation HMT ``feuillet surfacique \(2\pi\), feuillet volumétrique \(4\pi\)'' est
un autre objet. La section~\ref{sec:bisagra-hojas-pi-phi2} formalise son noyau
topologique : un tour de la base ferme le bord projeté en \(2\pi\) et
se termine sur le feuillet conjugué, tandis que le carré de la boucle ferme le revêtement
orienté en \(4\pi\). Le théorème est relatif au relèvement, à l'holonomie et au
lecteur géométrique déclarés. Son identification supplémentaire à une
représentation spinorielle requiert toujours le foncteur physique et ne se déduit pas
du partage des mêmes facteurs.

\subsection{Application entre multisections internes et représentations physiques}

L'objet particule, la projection
extension--voie, le sélecteur interne des 81 cellules, le centre électronique,
les treize multisections finies de famille, leur relèvement relatif au lecteur
cellulaire \(\pi_{81}\), le tritique de couleur, le registre, la signature, le caractère et
les tours forment la composition
\[
\text{extension}\longrightarrow\text{route}\longrightarrow
\text{multisection}\longrightarrow\text{signature}\longrightarrow
\text{caractère}\longrightarrow\text{tours}.
\]

L'obligation restante consiste à construire un foncteur de
réalisation qui apparie certaines étiquettes physiques individualisées et
composées aux multisections enrichies pertinentes et reproduise les
signatures raffinées déclarées sans consulter les masses ni les signatures cibles. Pour les
états composés, le registre est composé avant d'être projeté sur cinq
coordonnées. Le domaine du foncteur conserve l'électron central, l'atlas des
familles et la loi interne de caractère.


\paragraph{Typage probatoire de la persistance, de l'atlas et des familles.}
La persistance de voie, les anti-sections, la lecture bidirectionnelle et le
lecteur temporel tritique fixent l'architecture de départ. Le système
inverse, le sélecteur interne et la séparation des types réalisent cette
architecture au moyen de morphismes déclarés. Les recensements de l'atlas et les
identités de \(C_M\) sur les mots sont satisfaits dans tout le système fini ;
leur relèvement aux sections et le ruban de Möbius sont démontrés avec une
structure de départ explicite. Les coefficients électroniques
\((-22,+15)\) se déduisent avec la même structure explicite. Si un
relèvement ouvre au préalable une différence métrologique ou une signature cible,
sa valeur est une reconstruction calibrée.

\endgroup

\Needspace{18\baselineskip}
\section{Ontologie matérielle, charge, spin, statistique, couleur et saveur}
La définition complète est complétée par les domaines de représentation, les conditions de conjugaison, les complétions CAR/CCR, les secteurs de jauge, la composition, la virtualité et les contrôles de positivité.

\begingroup
\let\chapter\section
\let\section\subsection
\let\subsection\subsubsection
\let\subsubsection\paragraph
\let\clearpage\relax
\section{Ontologie HMT--MD des particules}
\Needspace{7\baselineskip}
\subsection{État généalogique matériel : extension, parcours, enregistrement, signature et représentation}
Un état HMT complet n'est pas un simple point. C'est le tuple

\[
 \mathfrak x=
 (\Gamma,\mathcal L_\Gamma,\sigma_\Gamma,\eta_\Gamma,
 \mathcal R_\Gamma,\mathcal B_\Gamma,\mathcal C_\Gamma),
\]
où \(\Gamma\) est un parcours survivant, \(\mathcal L_\Gamma\) son enregistrement généalogique, \(\sigma_\Gamma\) sa signature entière, \(\eta_\Gamma\) son raffinement harmonique, \(\mathcal R_\Gamma\) sa représentation, \(\mathcal B_\Gamma\) sa mémoire de retour et \(\mathcal C_\Gamma\) sa donnée de composition ou de frontière. Deux états ayant le même mot visible peuvent être distincts s'ils diffèrent par le préfixe, le bloc, la mémoire, la retenue, le feuillet ou l'ombre de Witt.

\Needspace{7\baselineskip}
\subsection{Particule comme section locale persistante de l'état généalogique dimensionnel}
Une particule est une classe d'équivalence de sections persistantes satisfaisant quatre conditions :

\begin{enumerate}
\item prolongement compatible à travers les cylindres du TPK ;
\item retour de phase avec mémoire d'état ;
\item représentation physique irréductible ou bloc invariant minimal ;
\item stabilité spectrale sous le couplage qui l'observe.
\end{enumerate}
Une section persistante est une famille compatible

\[
 s=(s_K)_{K\ge K_0},\qquad
 p_{L,K}(s_L)=s_K\quad(L\ge K\ge K_0).
\]
Deux sections \(s\) et \(s'\) représentent la même particule lorsqu'il existe un niveau \(K_1\) et une famille compatible de symétries internes \(g_K\in G_K\) tels que

\[
 s'_K=g_Ks_K\quad(K\ge K_1),
 \qquad
 p_{L,K}g_L=g_Kp_{L,K},
\]
et que les caractères observables de l'enregistrement généalogique coïncident. Cette relation identifie les descriptions qui diffèrent par la jauge interne, mais n'identifie pas les parcours de mémoire, retenue, feuillet, bloc ou holonomie distincts. La particule est la classe

\[
 [s]_{\rm pers}
 =\{s':s'\sim s\},
\]
et non un mot visible isolé ou une cellule choisie après connaissance de sa masse.

Si \(\widetilde\Gamma_K\) est l'extension au niveau \(K\), la persistance s'exprime par

\[
 p_{K+9,K}(\widetilde\Gamma_{K+9})
 =\widetilde\Gamma_K,\qquad
 g(K+9)=g(K),
\]
sans exiger

\[
 \mathcal B_{K+9}=\mathcal B_K.
\]
La particule n'est pas la phase qui revient, mais l'identité généalogique conservée à travers le changement d'état. Sa masse mesure la rigidité spectrale de cette persistance lorsqu'elle se couple à une réalisation matérielle.

\Needspace{7\baselineskip}
\subsection{Antiparticule comme image conjuguée d'une section persistante sous inversion de feuillet}
La conjugaison dodécaphasique est

\[
 C_M(\mathbf t)=-\operatorname{rev}(\mathbf t).
\]
La paire particule--antiparticule est formée de deux orientations d'une même bande généalogique. Si \(\chi\) désigne l'orientation de feuillet,

\[
 (\mathbf t,\chi)\sim(C_M\mathbf t,-\chi).
\]
Pour une énergie de bande

\[
 E_{\rm band}(\mathbf t,\chi)
 =E_+(\mathbf t)+\chi E_-(\mathbf t),
\]
avec

\[
 E_\pm(\mathbf t)
 =\frac12\bigl(E_0(\mathbf t)\pm E_0(C_M\mathbf t)\bigr),
\]
on obtient

\[
 E_{\rm band}(C_M\mathbf t,-\chi)
 =E_{\rm band}(\mathbf t,\chi).
\]
L'égalité des masses est donc une conséquence de la parité de l'opérateur sous conjugaison, tandis que charge et orientation changent. L'égalité des largeurs exige en outre que l'opérateur de désintégration soit conjugué par \(C_M\).

\Needspace{7\baselineskip}
\subsection{Masse comme caractère spectral de persistance de \(1\) parcours matériel}
La masse est la mesure spectrale, dans la droite \(\mathcal L_M\), de l'opérateur résultant de la compression du caractère, des tours, du lecteur bidirectionnel et de l'horloge sur le sous-espace persistant sélectionné par le couplage :

\[
 \widehat M_{P,a}
 =P_{P,a}\,
 \mathbf m_\star^{\rm ret}
 R_{\rm act}^{\widehat K_{P,a}/2}
 \widehat{\mathcal A}_{P,a}^{(0)}
 R_{\rm act}^{\widehat K_{P,a}/2}
 P_{P,a}.
\]
Le résultat général est

\[
 \mu_{P,a}^{\rm mass}(B)
 =\operatorname{Tr}\!\left(
 \rho_{P,a}E_{\widehat M_{P,a}}(B)\right).
\]
Il existe une masse scalaire \(m_P\) si et seulement si

\[
 P_{\operatorname{supp}\rho}\widehat M_{P,a}
 P_{\operatorname{supp}\rho}
 =m_P P_{\operatorname{supp}\rho},
\]
de manière équivalente, si la variance de \(\widehat M_{P,a}\) dans l'état préparé est nulle. Cette définition inclut l'électron de rang un, les multiplets imposés par Schur et les pôles isolés des résonances, sans les confondre.

\Needspace{7\baselineskip}
\subsection{Charge comme caractère de jauge transporté par la section matérielle}
La charge est le caractère orienté de la couche impaire de l'enregistrement généalogique sous conjugaison. Soit \(\sigma_{\rm odd}\) la signature impaire. Alors

\[
 Q(C_M\Gamma)=-Q(\Gamma),\qquad
 Q(\Gamma)=\mathfrak q\!\left(\sigma_{\rm odd}(\Gamma)\right),
\]
où \(\mathfrak q\) est l'homomorphisme de la signature impaire vers le réseau de charge de la réalisation. La signature massique paire et la signature de charge impaire sont des objets distincts : un mot peut avoir une somme linéaire nulle tout en portant une signature massique non nulle.

\Needspace{7\baselineskip}
\subsection{Spin comme représentation de l'action des rotations sur la fibre matérielle}
Le spin est la classe d'holonomie avec laquelle une section persistante répond à une rotation complète du repère d'orientation. La racine interne

\[
 J^2=-I
\]
sépare les deux clôtures :

\[
 U(2\pi)=
 \begin{cases}
 +I,&\text{secteur tensoriel},\\
 -I,&\text{secteur spinoriel},
 \end{cases}
 \qquad
 U(4\pi)=I.
\]
Un boson de spin entier appartient à une représentation qui se clôt à \(2\pi\) ; un fermion de spin demi-entier nécessite le double tour \(4\pi\). La valeur du spin ne se déduit pas d'un dénombrement cardinal isolé : c'est le poids de la représentation du groupe d'holonomie sur le sous-espace persistant. En particulier :

\[
 s=0:\text{ section scalaire},\qquad
 s=\tfrac12:\text{ bande spinorielle},
\]
\[
 s=1:\text{ section vectorielle transverse},\qquad
 s=2:\text{ composante tensorielle d'une réalisation métrique}.
\]
La dernière ligne définit un type de représentation ; elle n'impose pas l'existence d'une particule élémentaire de spin deux.

La chiralité et l'hélicité sont des données ultérieures et distinctes. Une involution orientée \(\Gamma_{\rm ch}^2=I\) définit

\[
 P_L=\frac{I-\Gamma_{\rm ch}}2,
 \qquad
 P_R=\frac{I+\Gamma_{\rm ch}}2,
\]
tandis que l'hélicité est le signe spectral du générateur de rotation suivant la direction de propagation. La première classe les feuillets de la représentation ; la seconde dépend de l'état cinématique. Cette séparation permet à un neutrino d'être fermionique par sa clôture \(4\pi\), même lorsque son couplage faible sélectionne une seule chiralité.

\Needspace{7\baselineskip}
\subsection{Statistique quantique induite par la symétrie d'échange des sections}
La statistique est la règle de complétion de l'espace des parcours identiques :

\[
 \mathcal F_-(\mathcal H)
 =\bigoplus_{n\ge0}\wedge^n\mathcal H,\qquad
 \mathcal F_+(\mathcal H)
 =\bigoplus_{n\ge0}\operatorname{Sym}^n\mathcal H.
\]
Dans la complétion extérieure, les relations CAR

\[
 \{a_i,a_j^\dagger\}=\delta_{ij},\qquad
 \{a_i,a_j\}=0
\]
impliquent

\[
 n_i^2=n_i
\]
et, par conséquent, l'exclusion de Pauli et la distribution de Fermi--Dirac. Dans la complétion symétrique, les relations CCR permettent une occupation multiple et conduisent à Bose--Einstein. La statistique ne s'infère pas de la masse : elle procède de la façon dont le parcours se compose sous échange.

\Needspace{7\baselineskip}
\subsection{Couleur comme représentation résiduelle de \(SU(3)\) sur les parcours de quarks}
La couleur est l'orbite résiduelle triadique d'une section de quark dans la fibre de charge et de saveur. Les trente-six positions colorées se décomposent comme

\[
 12_R+12_G+12_B.
\]
L'objet physique n'est pas l'une de ces couleurs, mais le triplet et son action :

\[
 \mathcal H_q^{\rm color}
 =\mathcal H_R\oplus\mathcal H_G\oplus\mathcal H_B.
\]
Si l'opérateur de masse commute avec l'action irréductible de couleur,

\[
 U_g\widehat M_qU_g^*=\widehat M_q,
\]
le lemme de Schur impose une masse commune dans le triplet. Choisir une cellule rouge, verte ou bleue comme représentante détruirait la covariance.

\Needspace{7\baselineskip}
\subsection{Saveur comme classe de parcours au sein d'une génération matérielle}
La saveur est une classe de couplage non équivalente au sein d'une même représentation de jauge, caractérisée par charge, génération, feuillet, signature fine, mémoire et projecteur d'interaction. Deux saveurs peuvent partager une fibre grossière et différer par l'entrelaceur qui les couple au lecteur.

Formellement, si \(F\) est la fibre commune, les saveurs \(f_i\) sont des représentations \(\pi_i\) pour lesquelles

\[
 \operatorname{Hom}_{G_F}(V_{f_i},\ell^2(F))\ne0
\]
et dont les projecteurs \(P_{f_i}\) ne sont pas conjugués par une symétrie de jauge qui demeure physique. La saveur n'est pas la profondeur de l'atlas et ne se réduit pas à un nom conventionnel.

\Needspace{7\baselineskip}
\subsection{Génération matérielle comme strate récurrente typée par profondeur de retour, signature fine et spectre du lecteur}
Une génération est une strate récurrente de la même grammaire de charge et de représentation, distinguée par profondeur de retour, signature fine et spectre du lecteur. L'étiquette de génération n'est physique que lorsqu'il existe un projecteur qui entrelace cette strate avec le canal de couplage. Une profondeur \(G_j\) ne doit donc pas être identifiée automatiquement à l'électron, au muon, au tau ou au neutrino stérile.

\Needspace{7\baselineskip}
\subsection{Champ de jauge comme connexion sur les fibres de charge}
Un champ de jauge est la réalisation d'une redondance de repère sur les fibres. Ses bosons élémentaires appartiennent à la direction nulle de l'opérateur de masse tant que la symétrie demeure exacte. Une masse de jauge ne s'obtient pas en approchant zéro par un caractère positif, mais par une brisure, un couplage longitudinal ou une réalisation de pôle déclarée.

\Needspace{7\baselineskip}
\subsection{Particule composée comme section persistante du produit de parcours liés}
Une particule composée est une section persistante de l'opérateur de composition, et non la somme des masses ou des signatures de ses constituants. Si \(\Gamma_1,\ldots,\Gamma_n\) sont les parcours constituants, on construit d'abord

\[
 \mathfrak C_{12}
 (\Gamma_1,\ldots,\Gamma_n;\mathcal T,\partial,\mathrm{carry})
 =\mathcal L_{\rm comp},
\]
où \(\mathcal T\) est l'arbre temporel et \(\partial\) la frontière de liaison. Ce n'est qu'ensuite qu'agit le lecteur

\[
 \mathcal L_{\rm comp}\longrightarrow\sigma_{\rm comp}
 \longrightarrow\chi_{\rm comp}
 \longrightarrow\widehat M_{\rm comp}.
\]
L'énergie de liaison réside dans le cocycle de composition, la mémoire de frontière et la modification du spectre ; elle n'est pas introduite comme correction externe.

\Needspace{7\baselineskip}
\subsection{Résonance comme section transitoire et largeur comme taux de désintégration}
Une résonance est un mode persistant non unitaire de l'opérateur couplé. Si

\[
 \lambda=re^{-i\theta},\qquad 0<r<1,
\]
alors

\[
 E=\frac{\hbar_{\rm int}}{t_0}
 \bigl(\theta-i\log r\bigr)
 =M-\frac{i}{2}\Gamma.
\]
La masse \(M\) et la largeur \(\Gamma\) sont les parties réelle et imaginaire d'un pôle, mais exigent des lecteurs distincts. La largeur peut aussi s'obtenir à partir d'un semi-groupe comprimé :

\[
 T_X=Q_XUQ_X,\qquad
 \lambda_X=-\frac1{t_0}\log\rho(T_X),\qquad
 \Gamma_X=\hbar_{\rm int}\lambda_X.
\]
On ne déduit pas une largeur du seul fait qu'un parcours possède un horizon fini.

\Needspace{7\baselineskip}
\subsection{Section virtuelle comme extension intermédiaire non asymptotique de l'atlas}
Une section est virtuelle lorsque son prolongement se termine de manière non triviale :

\[
 \operatorname{Ext}_{L}(s)\ne\varnothing,\qquad
 \operatorname{Ext}_{L+1}(s)=\varnothing,\qquad
 \mathcal L(s)=L<\infty.
\]
Son contenu primaire est l'horizon \(L\), l'obstruction terminale, l'enregistrement généalogique fini et l'amplitude de propagation. Elle n'est ni une particule de masse imaginaire ni un état documentaire non encore classé.

\Needspace{7\baselineskip}

\endgroup
\begingroup
\let\chapter\section
\let\section\subsection
\let\subsection\subsubsection
\let\subsubsection\paragraph
\let\clearpage\relax
\section{Classification complète des particules et des états physiques}
\label{sec:catalogo-realizaciones-tipadas-rev8}
\index{secteurs physiques!inventaire typé}
\index{atlas!81/56/13/324}

Les réalisations sont ordonnées par objet, multiplicité et codomaine. Le domaine mathématique contient \(81\) positions, regroupées en
codomaines distincts. Le domaine mathématique contient \(81\) positions, regroupées en
\(56\) familles de trajectoires et \(13\) classes d'extension ; les \(324\)
trajectoires orientées ajoutent l'orientation initiale et ne représentent pas de masses
ou de particules supplémentaires. La colonne ``opérateur d'évaluation'' n'impose pas de masse là où l'objet exige
une carte nulle, une composition préalable, un invariant topologique ou une
obstruction au prolongement.

\subsection{Définitions spectrales de particule, spin, couleur et saveur}

Une \emph{particule} est une section survivante réalisée comme parcours
orienté, munie d'un enregistrement, d'une représentation de symétrie, d'un opérateur spectral
et d'une carte physique. Un fermion réalise la graduation extérieure et les relations
CAR ; un boson réalise la complétion symétrique et les relations CCR. Le
spin enregistre l'holonomie de l'orientation : le retour de \(2\pi\) clôt
la phase bosonique, tandis que la section spinorielle retrouve son état après
\(4\pi\). Cette distinction relie le double feuillet surfacique--volumétrique à la
parité et à l'exclusion de Pauli.

La \emph{couleur} est la représentation ternaire qui relève la coordonnée
\(r-c\pmod3\) vers l'orbite fondamentale puis vers l'adjointe de
\(\mathfrak{su}_3\). La \emph{saveur} identifie une base de réalisation
au sein d'une représentation de charge, couleur et spin fixés ; CKM et PMNS
assurent le transport entre la base d'interaction et la base spectrale. Les neutrinos
sont des sections fermioniques neutres : leurs étiquettes de saveur relèvent du
lecteur de mélange et leurs masses des vecteurs propres de l'opérateur correspondant.

Une particule virtuelle est une section finie avec enregistrement et propagateur jusqu'à
un horizon de prolongement ; une réalisation de Majorana est le secteur
fermionique réel fixé par conjugaison ; et les secteurs de Fibonacci, d'Ising et de
\(\mathbb Z/6\mathbb Z\) donnent des dimensions quantiques, des règles de fusion et
des caractères de tresse. Chaque définition conserve ainsi le codomaine que la
structure lui assigne.


\endgroup
\begingroup
\let\chapter\section
\let\section\subsection
\let\subsection\subsubsection
\let\subsubsection\paragraph
\let\clearpage\relax

Le tableau conserve en outre les réalisations quantiques que sa classification
abrège. Les trois leptons chargés et les six saveurs de quarks appartiennent à la
complétion extérieure : CAR conduit à Pauli et à Fermi--Dirac, sans attribuer ces
relations à un TRIT isolé. Les quarks occupent trente-six cellules lorsque l'on
conserve couleur et antiparticule avant le quotient. Le support neutre commun
\(\mathcal H_\nu=\bigoplus_{j=1}^{4}\mathcal H_{\nu,G_j}\) a pour dimensions
\(2+4+6+8=20\). Les symboles \(G_j\) désignent des strates de profondeur de l'atlas ;
\(f_e,f_\mu,f_\tau\) désignent la base d'interaction. En particulier, \(G_4\) est
une strate du support neutre commun et ne s'identifie pas à un état stérile ou
exotique sans projecteur supplémentaire. PMNS transporte la base d'interaction vers la
base des vecteurs propres ; il ne crée ni valeurs propres, ni échelle de masse, ni sélecteur de
profondeur. Le photon conserve aller et retour dans la
carte de jauge nulle : \(m=0\) ne signifie pas évaluer \(\chi_0\) avec un exposant très
négatif. \(W^\pm\) et \(Z^0\) sont des représentations bosoniques massives, avec
conjugaison de charge pour \(W^+\leftrightarrow W^-\). Enfin,
l'antisection s'obtient par inversion et renversement de bande ; la masse est paire
sous cette involution et n'est pas doublée comme second paramètre indépendant.

Tous les secteurs conservent le préfixe causal
\[
\text{extension survivante}\longrightarrow\text{route observable}
\longrightarrow\text{registre chronologique enrichi}.
\]
Le lecteur change ensuite. Dans le secteur matériel positif, la chaîne se poursuit par
\(\text{signature}\to\text{caractère}\to\text{tours}\to\text{carte dimensionnelle}\).
Le secteur nul évite le caractère positif et fixe \(m=0\) ; les hadrons composent
d'abord des enregistrements chronologiques enrichis ; les secteurs topologiques utilisent des invariants de
fusion ou de tresse ; et la virtualité se reconnaît à l'obstruction terminale. Le
domaine de ces applications est la classification complète précédente. Un essai
numérique sur un sous-ensemble n'établit que le cas fini testé et ne
réduit pas ce domaine.

Les dénombrements qui accompagnent ce tableau relèvent de quotients différents :
\[
81=25_{\ell^-}+20_\nu+20_d+16_u,
\]
\[
56\ \text{familles à fibres}\quad
41\times1+10\times2+2\times3+1\times4+2\times5=81,
\]
\[
13\ \text{multisections},\qquad324=81\times4\ \text{généalogies}.
\]
De même, les \(192\)-A classes de représentation associées à
\(\operatorname{diag}(2,6,2,2,4)\) et l'indice réticulaire \(192\)-B de
\(\operatorname{SNF}(2,2,2,2,12)\), avec \(52\) signatures brutes, ne sont pas le
Le cardinal du catalogue historique d'exposition ne compte pas les particules individuelles et ne définit pas la taxonomie matérielle.
non des particules individuelles.

\endgroup

\begingroup
\let\chapter\section
\let\section\subsection
\let\subsection\subsubsection
\let\subsubsection\paragraph
\let\clearpage\relax
% Vista científica exacta materializada desde una rama condicional.
% Fuente: source/demostraciones_primarias/27_04_particula_estadistica.tex; rama: HMTIncludeParticleCore.
% No modifica el contenido matemático de la rama de origen.
\chapter{Sections persistantes et propagation finie}
\label{chap:particula}

La notion mathématique de particule est introduite après la structure de frontière et avant la spécialisation de sa dynamique. Son support primaire est une extension survivante du système d’états TPK\@. Une trajectoire orientée est l’ombre séquentielle de cette extension ; sur elle, les degrés internes forment un fibré discret. Le propagateur, la somme sur les chemins et les représentations de Clifford sont introduits après la construction de la droite d’action. Cette hiérarchie empêche qu’une particule soit définie par le choix externe d’une ligne ou d’un parcours et distingue l’objet mathématique de son identification ultérieure à une espèce physique.

\section{Extension survivante et projection observable}

Soit \((\mathcal S_m,\rho_{m+1,m})\) le système inverse des états qui franchissent les clôtures jusqu’à la profondeur \(m\). Un état global est une famille compatible
\[
 x=(x_m)_m\in\Omega_{\HMT}^{\rm surv}
 :=\varprojlim_m\mathcal S_m.
\]
Une carte séquentielle de profondeur \(R\) produit un parcours
\[
 \gamma_x^{[R]}=\operatorname{sh}_R(x)
 =(c_0\longrightarrow\cdots\longrightarrow c_R).
\]
Le symbole \(\operatorname{sh}\) désigne une lecture d’ombre : différentes cartes peuvent présenter un même état global, pourvu qu’elles conservent le feuillet, l’orientation, le bord et l’enregistrement. Dans le \cref{chap:extension-ruta-caracter} est construit l’enregistrement enrichi qui fait descendre le caractère d’action de ces présentations à l’extension.

\section{Particule comme section compatible d’un fibré discret}

Soit \(\Gamma_{\rm TPK}\) le graphe dont les sommets sont les états complets du noyau de phase topologique et dont les arêtes sont les transformations admissibles. Pour \(x\in\Omega_{\HMT}^{\rm surv}\) et sa trajectoire observée
\[
 \gamma_x=(c_0\longrightarrow c_1\longrightarrow\cdots
 \longrightarrow c_R),
\]
soit \(\operatorname{Typ}(c_t)\) l’ensemble fini des types arithmétiques, des orientations et des mémoires sur l’état \(c_t\), et soit \(E_{c_t}=\mathbb C[\operatorname{Typ}(c_t)]\) l’espace vectoriel complexe libre engendré par ces types. Le fibré discret restreint à \(\gamma_x\) est
\[
 \mathcal B_{\gamma_x}=\bigsqcup_{t=0}^{R}E_{c_t}
 \longrightarrow\{0,\ldots,R\}.
\]
Chaque arête \(c_t\to c_{t+1}\) induit une application linéaire de transport \(\tau_t:E_{c_t}\to E_{c_{t+1}}\).

\begin{definition}[Particule mathématique]
Une particule de Mécanique Dimensionnelle est un tuple
\[
 \mathfrak p=(x,\gamma_x,\mathcal B_{\gamma_x},s,n,U),
\]
où \(x\in\Omega_{\HMT}^{\rm surv}\) est l’extension survivante et \(\gamma_x=\operatorname{sh}_R(x)\) son ombre dans la carte déclarée ; de plus,
\[
 s:\{0,\ldots,R\}\longrightarrow\mathcal B_{\gamma_x},\qquad
 s(t)\in E_{c_t},\qquad s(t+1)=\tau_t(s(t)),
\]
est une section compatible ;
\[
 n:\bigsqcup_{t=0}^{R}\operatorname{Modes}(E_{c_t})
 \longrightarrow\mathbb N
\]
est sa fonction d’occupation, où \(\operatorname{Modes}(E_{c_t}):=\operatorname{Typ}(c_t)\) désigne la base distinguée de l’espace vectoriel libre ; et
\[
 U_{t_2,t_1}:E_{c_{t_1}}\longrightarrow E_{c_{t_2}}
\]
est une famille de propagateurs linéaires qui entrelace les transports, \(U_{t,t}=\operatorname{id}_{E_{c_t}}\), \(U_{t+1,t}=\tau_t\), et vérifie
\[
 U_{t_3,t_2}\circ U_{t_2,t_1}=U_{t_3,t_1}
 \qquad(0\leq t_1\leq t_2\leq t_3\leq R).
\]
Chaque composition est définie entre les fibres correspondantes. L’identification de \(\mathfrak p\) à une espèce physique requiert en outre une représentation du groupe de symétrie pertinent et une application des états HMT vers les états de cette représentation.
\end{definition}

L’extension conserve la généalogie et l’enregistrement ; le parcours en expose la lecture ordonnée ; les degrés internes appartiennent aux fibres, l’occupation à \(n\) et l’évolution à \(U\). La particule est la compatibilité de ces niveaux, et non l’un d’eux pris séparément. En particulier, le parcours n’est pas choisi par la particule : il est une projection de l’extension qui la réalise.

\endgroup
\begingroup
\let\chapter\section
\let\section\subsection
\let\subsection\subsubsection
\let\subsubsection\paragraph
\let\clearpage\relax
% Fragmento científico exclusivo de 01_electron.tex, líneas 120--197.
% El prefijo 1--118 coincide byte a byte con la ontología principal ya
% publicada; este fragmento conserva únicamente la formulación cuántica
% adicional sin repetir definiciones idénticas.
\section{Ontologie quantique de la particule}
\Needspace{7\baselineskip}
\subsection{État, histoire et réalisation}
Un état généalogique est une section compatible du système d'extensions du TPK. Une histoire observable est la projection chronologique de cette section sur un lecteur. Une réalisation physique est la représentation de l'histoire dans une droite dimensionnelle, un espace de Hilbert, une algèbre de champs ou un codomaine topologique. Ces trois notions ne sont pas interchangeables. L'état porte l'identité ; l'histoire porte la preuve ; la réalisation porte l'observable.

Soit \(\mathcal X_{\rm TPK}^{\rm enr}\) l'espace des états enrichis et \(p\) le lecteur observable. Une particule est une classe \([X]\) de sections persistantes telle que

\[
p(X_{K+1})|_K=p(X_K),
\]
la loi de transition conserve les relations de l'enregistrement généalogique et que le prolongement ne se termine pas sur un horizon d'obstruction. Le parcours

\[
\gamma_X=(p(X_0),p(X_1),p(X_2),\ldots)
\]
est l'histoire visible de la particule. Il n'est pas, par définition, une trajectoire spatiale. Une réalisation métrique ultérieure peut le convertir en courbe, orbite ou propagateur lorsqu'elle construit l'entrelaceur correspondant.

L'identité physique se conserve sous raffinement lorsque les carrés

\[
\begin{CD}
X_{K+1} @>{F_{K+1}}>> X'_{K+1}\\
@V{r_{K+1,K}}VV @VV{r'_{K+1,K}}V\\
X_K @>>{F_K}> X'_K
\end{CD}
\]
commutent. La masse, la charge et le spin sont des lecteurs distincts de ce même antécédent. L'égalité de l'un n'impose pas l'égalité des autres.

\Needspace{7\baselineskip}
\subsection{Représentation quantique de la particule comme section locale persistante de l'état généalogique}
Une particule HMT--MD est une section persistante d'une extension enrichie qui satisfait quatre conditions :

\begin{enumerate}
\item elle possède un parcours observable compatible avec la monodromie nonadique ;
\item elle conserve un enregistrement généalogique sous raffinement ;
\item elle admet une représentation d'orientation, de charge, de spin et de statistique ;
\item elle présente un résultat spectral stable dans le codomaine physique déclaré.
\end{enumerate}
La stabilité ne signifie pas l'immobilité. Elle signifie que les changements de phase, de feuillet et de mémoire appartiennent à une même classe généalogique. La particule peut clore sa phase sans clore son état, changer de feuillet tout en conservant sa masse, ou traverser plusieurs cellules en maintenant la même identité d'extension.

Une espèce est une famille de particules partageant le classificateur interne

\[
\mathfrak s=
(\mathfrak q,\mathfrak j,\mathfrak c,\mathfrak f,\mathfrak g,
\mathfrak r,\mathfrak o),
\]
où \(\mathfrak q\) est la représentation de charge, \(\mathfrak j\) l'holonomie de spin, \(\mathfrak c\) l'orbite de couleur, \(\mathfrak f\) la saveur, \(\mathfrak g\) la génération, \(\mathfrak r\) la famille de parcours et \(\mathfrak o\) l'opérateur de masse. Une espèce non centrale peut correspondre à une multisection complète ; exiger d'elle un point APP unique avant de déclarer le couplage physique constitue une erreur de type.

\Needspace{7\baselineskip}
\subsection{Conjugaison quantique de l'antiparticule par inversion de feuillet d'une section persistante}
L'involution dodécaphasique est

\[
C_M(\tau)=-\operatorname{rev}(\tau).
\]
Elle agit sur l'histoire orientée et se prolonge à l'enregistrement généalogique. La charge et l'orientation linéaire sont impaires ; la mémoire quadratique est paire. Si \(q\) est le lecteur de charge,

\[
q(C_MX)=-q(X).
\]
L'égalité des masses exige en outre la covariance opératorielle. Si \(K\) réalise la conjugaison dans l'espace physique et

\[
K\widehat M_XK^{-1}=\widehat M_{\bar X},
\]
en transportant aussi les canaux et la prescription causale lorsque \(K\) est antiunitaire, alors \(m_{\bar X}=m_X\). La seule parité de l'enregistrement ne remplace pas cette prémisse.

Particule et antiparticule sont deux sections orientées d'une même bande généalogique. Dans la réalisation de Möbius, un tour échange le feuillet et deux tours rétablissent l'orientation. La conclusion ne procède pas du dessin d'un ruban tordu : elle exige l'involution, le relèvement au revêtement double, la boucle de retour et l'holonomie de signe.

Pour l'enregistrement électronique \(\tau_e=+++---+++---\),

\[
C_M\tau_e=\tau_e.
\]
Le mot central est fixe, tandis que la représentation de charge distingue électron et positron. On obtient ainsi une même masse sans doubler le parcours temporel. La clôture à 216 pas ne signifie pas 108 pas de particule et 108 d'antiparticule : elle signifie le rétablissement du relèvement cinématique orienté.

\Needspace{7\baselineskip}

\endgroup
\begingroup
\let\chapter\section
\let\section\subsection
\let\subsection\subsubsection
\let\subsubsection\paragraph
\let\clearpage\relax
La Loi Générale des Masses est formulée comme une transformation généalogique et
spectrale, et non comme une collection de nombres. Son domaine effectif est la chaîne
\[
\text{APP}\longrightarrow\text{TRIT}\longrightarrow\text{TPK}
\longrightarrow\Gamma^{\rm surv}\longrightarrow\Gamma^{\rm obs}
\longrightarrow L_{108}\longrightarrow\sigma\in\mathbb Z^5
\longrightarrow\chi_{\rm full}\longrightarrow\widehat M,
\]
suivie, lorsque la représentation physique le permet, de la fonctionnelle qui
extrait une coordonnée scalaire de la droite de masse. Chaque flèche conserve son
domaine, son orientation, sa mémoire et son statut ; aucune valeur externe ne choisit
rétroactivement un parcours, une signature ou un coefficient harmonique.

L'atlas fini contient 81 cellules, 56 familles de parcours, 13 multisections et
324 parcours orientés. Le point central électronique est l'unique fibre ponctuelle
équivariante. Les autres classes élémentaires se réalisent naturellement comme
fibres, orbites ou multisections, si bien que leur résultat primaire peut être un
opérateur ou un spectre avant d'être une masse scalaire. Cette différence n'est pas une
lacune du domaine, mais une propriété de la géométrie interne.

Dans la fibre centrale, le sélecteur de base et le lecteur bidirectionnel produisent
\[
\mathfrak e_0
=\frac{\sqrt3}{4}
 \left[\exp\!\left(\frac{\varphi}{\pi^2}\right)
       -22\alpha_{\rm HMT}^3\right](1+15\Delta_4),
\qquad
\kappa_e=80-54\alpha_{\rm HMT}+6\Delta_4,
\]
\[
\mathfrak e_e^\beta=\mathfrak e_0R_{\rm act}^{\kappa_e},
\qquad
\varsigma_{u_E}(\mathfrak e_e^\beta)
=0.510998950690000808608\ldots\ \mathrm{MeV}.
\]
L'échelle absolue d'action demeure dans la branche
\(\Sigma_H=\varphi\alpha_{\rm HMT}^{16}\), d'ordre décimal \(-34\), et se
couple à l'horloge par \(\hbar_{\rm int}\omega\) ; dans les quotients de
masses, toute section commune se simplifie et seuls subsistent les caractères, les
lecteurs et les raffinements relatifs correctement typés.

La hiérarchie harmonique ne consiste pas en huit corrections isolées. C'est le semi-groupe
global
\[
\langle90,120\rangle
=\{90,120\}\cup\{30r:r\ge6\},
\qquad
c_n=q_-^n+q_+^n,
\quad s_n=q_-^n-q_+^n,
\]
agissant sur un unique opérateur bicouche. Les niveaux visibles sont des témoins de
cette hiérarchie ; sa complétude représentative ne remplace pas l'application physique qui
doit transporter parcours, retenue, mémoire et enroulement vers les coefficients de chaque
état.

Les réalisations physiques typées, les dix-sept évaluations conditionnelles, les
deux sections nulles et les secteurs neutriniques, composés, résonants,
topologiques et virtuels sont ainsi réunis sans aplanir leurs codomaines.
L'inventaire ultérieur de 471 masses et de 384 largeurs conserve la confrontation
état par état, mais ne se confond pas avec un dénombrement de prédictions. Une
comparaison quantitative n'acquiert une force physique qu'après déclaration de
l'observable, de l'unité, du schéma, de l'échelle, de l'incertitude, de la covariance et de la règle de
réalisation.

La conclusion structurelle est précise : la masse est la mesure spectrale d'une
persistance sous couplage. La particule est la section persistante ; le parcours,
son histoire observable ; la signature, sa lecture entière ; le caractère, son action
adimensionnelle ; les tours, sa résolution harmonique ; et la carte dimensionnelle, la
réalisation de cette généalogie en unités physiques. L'atlas et les inventaires
précédents permettent de reproduire chaque affirmation sans réduire la théorie à un
tableau abrégé ni transformer une interface ouverte en théorème établi.

\endgroup

\section{Composition et identité}
Une particule composée conserve son identité lorsque le registre de liaison, la frontière, la représentation et les observables demeurent compatibles sous raffinement. Une résonance ajoute un pôle et une largeur ; elle ne constitue pas une nouvelle coordonnée de la signature de masse. Le critère d’identité est réfuté s’il dépend de la valeur externe ou si deux histoires inéquivalentes sont identifiées sans isomorphisme de sections.
